% docs/manual/developer/chapters/08-runner-cli.tex
% 第 8 章：runner & CLI

\chapter{runner 与 CLI 子系统}

\modname{runner} 是 \openbench{} 的"编排"层：把已校验的 config 翻译为 task 列表、调度执行、聚合结果。\modname{cli} 是用户接触的最外层。本章把两者一起讲。

\section{runner.local：主流程}

\file{runner/local.py}（906 行）是主入口 \texttt{run\_evaluation(cfg)}：

\begin{enumerate}
  \item 设 \texttt{HDF5\_USE\_FILE\_LOCKING=FALSE}（避免多 worker 锁定 NC）
  \item 调 \modname{config.adapter.build\_runner\_bindings(cfg)} → 拿 RunnerConfig + namelists
  \item 建输出目录 \file{output/<run>/\{data,metrics,scores,figures,comparisons,reports,scratch,tmp\}/}
  \item 构建 task 列表（var × sim × ref 三元组）
  \item Phase 1: Evaluation
  \item Phase 2: Comparison（如启用）
  \item Phase 3: Groupby（如启用）
  \item Phase 4: Statistics（如启用）
  \item Phase 5: Report（如启用）
\end{enumerate}

\subsection{task 数据结构}

每个 task 是一个 dict：

\begin{minted}{python}
task = {
    "var_name": "Evapotranspiration",
    "sim_source": "CoLM2024_test",
    "ref_source": "GLEAM_v4.2a_LowRes",
    "bindings": <Bindings object>,
    "cache_key": "...",
    "config_hash": "...",
    "use_cache": True,
    "cache_dir": "output/<run>",
}
\end{minted}

\subsection{增量缓存命中流程}

\begin{enumerate}
  \item Per-task：算 \texttt{config\_hash}（覆盖 var/sim/ref/metrics/scores 等）
  \item 查 \file{output/<run>/scratch/cache/<cache\_key>.json}
  \item Hash 一致 → 通过 \modname{\_find\_existing\_outputs} 检查输出文件确实存在
  \item Hash 不一致或输出缺失 → 重新计算 + 写新 cache
\end{enumerate}

\subsection{cache 失效的精确度（多 ref 关键）}

\texttt{cache\_key = make\_cache\_key(var, sim, ref)} 是 \emph{三元组}。\modname{\_find\_existing\_outputs} 用精确前缀 glob：

\begin{minted}{python}
patterns = [
    f"{glob.escape(var)}_ref_{glob.escape(ref)}_sim_{glob.escape(sim)}_*",
    f"{glob.escape(var)}_stn_{glob.escape(ref)}_{glob.escape(sim)}_evaluations*",
]
\end{minted}

历史问题：早期 pattern 是 \texttt{\{var\}\_*\{ref\}*\{sim\}*}，子串通配会让 \texttt{RefA} 误匹配 \texttt{RefAB} 的输出。修复后多 ref 配置下不同 ref 的 cache 失效彼此独立。

\subsection{\_preprocess\_variable：dedupe 与 unified\_mask 备份}

\modname{runner.local.\_preprocess\_variable}（约 200 行）是多 ref / 多 sim 正确性的核心。每变量调用一次，按以下规则跑预处理：

\begin{itemize}
  \item \textbf{Dedupe key} = \texttt{(ref\_source, "\_grid" if both grid else sim\_source)}
  \begin{itemize}
    \item Grid×grid：跨 sim 共享 ref 预处理（写一份 flat NC）
    \item 任一边 stn：每对独立，因为 \modname{processing.extract\_station\_data\_if\_needed} 会 \emph{删除} flat NC
  \end{itemize}
  \item \textbf{Flat NC backup-restore}（保护 unified\_mask 累积态）：
  \begin{itemize}
    \item Grid prep 完成 / unified\_mask 写回后立即 \texttt{shutil.copy2} 到 \texttt{<flat>.unifiedmask.bak}
    \item Stn prep 之前再次备份（防止删除前没存）
    \item 后续 grid 任务发现 flat 缺失 → 从 .bak 恢复（保留累积 mask）
    \item 循环末尾：若仍有缺失，从 .bak 恢复；最终清理 .bak
  \end{itemize}
  \item \textbf{Stn×stn 优化}：同 ref + 第二个 stn sim 时，符号链接 ref 文件而非重新抽样
\end{itemize}

历史 bug：早期 \texttt{refs\_done: bool} 让多 ref 中只第一个被预处理（\texttt{c9fcbdc} 后修为 \texttt{set[str]}）；\texttt{has\_grid\_evaluation} 仅检查 \texttt{sim\_sources[0]} 让混合类型评估错跳过 phase（\texttt{48bee6f} 改完整笛卡尔积）。

\subsection{评估阶段任务循环}

\modname{run\_evaluation} 主循环：

\begin{enumerate}
  \item 调 \modname{\_preprocess\_variable} 处理每个变量的全部 (sim, ref)
  \item 对每个 task 调 \modname{\_evaluate\_single}：
  \begin{enumerate}
    \item Cache 检查（hash + 输出文件存在）→ 命中即返回 \texttt{skipped=True}
    \item 创建 \modname{Evaluation\_grid} 或 \modname{Evaluation\_stn} 实例
    \item 调 \texttt{make\_Evaluation()} / \texttt{make\_evaluation\_P()} 运行
    \item 写新 cache
  \end{enumerate}
  \item 按 \yamlkey{num\_cores} 决定串行 or 并行（当前 \modname{runner} 层未启用变量级并行；并行在 \modname{DatasetProcessing} 内部）
\end{enumerate}

\subsection{错误传播}

\begin{itemize}
  \item 单 task 失败 → 记录到 \texttt{results["errors"]}，继续其他 task
  \item 全部失败 → \texttt{status=failure}，run\_evaluation 抛异常
  \item 部分失败 → \texttt{status=partial}，run\_evaluation 正常返回但 errors 非空
  \item CLI \cli{run} 据此决定退出码
\end{itemize}

\section{runner.cache：缓存实现}

\file{runner/cache.py}（146 行）：

\begin{itemize}
  \item \texttt{make\_cache\_key(var, sim, ref)} → 构造确定的 cache key 字符串
  \item \texttt{EvaluationCache} 类：load/save JSON cache
  \item \texttt{EvaluationCache.hash\_config(d)} → SHA-256 of canonical JSON
\end{itemize}

\noteinline{cache 写入是 atomic（先写 \texttt{.tmp} 再 rename），避免崩溃时损坏。}

\section{runner.remote：SSH 现状}

\file{runner/remote.py}（651 行）含完整 SSH 框架：连接、文件 sync、远程执行、日志拉取。但当前 v3.0a1 \emph{未在 CLI 暴露}：

\begin{itemize}
  \item \cli{openbench run --remote HOST} 立即 SystemExit(1) 并提示用 GUI
  \item GUI 的 Remote Run 入口实际调用此模块
  \item 后续版本会把 CLI 远程入口接通
\end{itemize}

详见运维卷第 3 章。

\section{CLI 主入口：cli.main}

\file{cli/main.py} 用 click + \texttt{LazyGroup} 自定义类：

\begin{minted}{python}
class LazyGroup(click.Group):
    COMMAND_MAP = {
        "run": "openbench.cli.run:run",
        "check": "openbench.cli.check:check",
        # ...
    }

    def get_command(self, ctx, cmd_name):
        if cmd_name in self.COMMAND_MAP:
            module_path, attr = self.COMMAND_MAP[cmd_name].rsplit(":", 1)
            mod = importlib.import_module(module_path)
            return getattr(mod, attr)
\end{minted}

\subsection{为什么 lazy load}

\openbench{} 的命令模块各自 import \modname{config}, \modname{data.registry}, \modname{runner} 等重模块。如果 \cli{openbench --help} 启动时全部加载，启动时间 > 1 秒。LazyGroup 让 \cli{--help} 只 import \modname{cli.main}（< 100ms）；具体子命令 import 推迟到首次执行。

\subsection{加新命令}

\begin{enumerate}
  \item 建 \file{cli/<newcmd>.py}，定义 \texttt{@click.command()} 装饰的函数
  \item 在 \texttt{LazyGroup.COMMAND\_MAP} 加映射
  \item 在 \texttt{list\_commands} 里加名字（让 \cli{--help} 显示）
\end{enumerate}

\section{Walkthrough：添加新 CLI 命令}

假设要加 \cli{openbench export-stats CONFIG -o stats.csv}（导出某次评估的所有 statistic 表合并为一个 CSV）：

\subsection{Step 1：建 \file{cli/export\_stats.py}}

\begin{minted}{python}
"""openbench export-stats command."""
import click
import pandas as pd
from pathlib import Path


@click.command("export-stats")
@click.argument("config", type=click.Path(exists=True))
@click.option("-o", "--output", default="stats.csv", help="Output CSV path.")
def export_stats(config, output):
    """Export all statistic tables into one CSV."""
    from openbench.config import load_config

    cfg = load_config(config)
    output_dir = Path(cfg.project.output_dir) / cfg.project.name / "statistics"

    if not output_dir.exists():
        click.secho(f"No statistics dir at {output_dir}; run evaluation first.", fg="red")
        raise SystemExit(1)

    dfs = []
    for csv_path in output_dir.rglob("*.csv"):
        df = pd.read_csv(csv_path)
        df["statistic"] = csv_path.parent.name
        dfs.append(df)

    if not dfs:
        click.secho("No statistic CSVs found.", fg="yellow")
        raise SystemExit(1)

    merged = pd.concat(dfs, ignore_index=True)
    merged.to_csv(output, index=False)
    click.secho(f"✓ Wrote {len(merged)} rows to {output}", fg="green")
\end{minted}

\subsection{Step 2：注册到 LazyGroup}

\file{cli/main.py}：

\begin{minted}{python}
class LazyGroup(click.Group):
    COMMAND_MAP = {
        # ... 现有
        "export-stats": "openbench.cli.export_stats:export_stats",
    }

    def list_commands(self, ctx):
        return [..., "export-stats"]
\end{minted}

\subsection{Step 3：测试}

\file{tests/test\_cli\_integration.py}：

\begin{minted}{python}
def test_export_stats(tmp_path):
    runner = CliRunner()
    # 准备一个有 statistics/ 输出的 mock 目录
    # ...
    result = runner.invoke(cli, ["export-stats", str(cfg), "-o", out])
    assert result.exit_code == 0
    assert out.exists()
\end{minted}

\subsection{Step 4：文档}

加到用户卷附录 B（CLI 参考）+ 第 6 章相关位置。

\section{cli 子模块约定}

\begin{itemize}
  \item 命令模块文件名：snake\_case 与命令名一致（\file{cli/data.py} → \cli{openbench data}）
  \item 参数：用 \texttt{@click.argument} 表示必填位置参数；\texttt{@click.option} 表示带名称的选项
  \item 退出码：成功 0；失败 \texttt{raise SystemExit(N)}（N 通常 1）
  \item 输出：用 \texttt{click.secho} 带颜色；错误用 \texttt{fg="red"}
  \item 子命令组：用 \texttt{@click.group()} 定义父命令（如 \cli{data}），子命令用 \texttt{@<group>.command()}
\end{itemize}

\section{常见运维问题}

\subsection{运行卡住但不报错}

通常是某个 task 在数据 IO 阶段挂起。看 \file{output/<run>/run.log} 最后几行，定位卡住的 task。

\subsection{HDF5 lock 错误}

\openbench{} 在 \cli{run} 启动时已设环境变量。如果在 Notebook / 脚本里手动调 \texttt{run\_evaluation}，需自己设：

\begin{minted}{python}
import os
os.environ["HDF5_USE_FILE_LOCKING"] = "FALSE"
from openbench.runner.local import run_evaluation
\end{minted}

\subsection{partial 退出但用户希望 fail-fast}

当前 \openbench{} 的策略是"宽容"——单 task 失败继续。如果你要 fail-fast，需要改 \modname{runner.local.run\_evaluation} 的错误传播逻辑。

\section{下一步}

下一章详解 \modname{gui} 扩展：14 个 wizard page 的架构与添加新 page 的 walkthrough。
